\scanpage{67}
\mspage{1}
Soit $K$ un sous-corps du corps des complexes $\mathbb C$ -- des cas intéressants seront celui où $K=\mathbb Q$, ou $K$ extension finie de $\mathbb Q$, ou $K=\widetilde{\mathbb Q}$ extension abélienne maximale de $\mathbb Q$ dans $\mathbb C$ -- je serai intéressé en tous cas avant tout \unsure{des|aux} cas où $K$ est de t.f. sur $\mathbb Q$. (Le cas de $\widetilde{\mathbb Q}$ est un cas ``limite'' de tels cas, pour ne pas être gêné par les racines de $1\ldots$).

On s'intéresse aux schémas \inserted{$X$} de t.f. sur $K$, et aux espaces \inserted{analytiques} $X(\mathbb C)$ correspondants. Comme on s'intéresse aux invariants ``topologiques'' de $X,X(\mathbb C)$, on traitera les morphismes $X'-X$ qui induisent des homéomorphismes universels comme des iso -- en particulier les éléments nilpotents ne comptent pas -- une imm. fermée surjective sera traitée comme un iso.

Un langage commode est de considérer $X$ comme donné par l'espace analytique (ou \unsure{$\overline{\mathrm m}$|$=$} seulement topologique) $X(\mathbb C)$, plus une \uline{structure} supplémentaire dessus : la structure de ``var. alg. sur $K$'', à coups de top. de Zariski, de faisceau des fonctions $K$-polynomiales etc. On écrira par la suite $X$ au lieu de $X(\mathbb C)$.

On distingue sur $X$ le groupoïde fondamental $\Pi_1(X)$ au sens topologique, et son compactifié profini $\widehat\Pi_1(X)$, qui \deleted{\unsure{est en un sens}} \inserted{peut} aussi être défini directement en termes du $\mathbb C$-schéma $X$. On a un foncteur ``d'inclusion''
\margin{66}
